The literature supports a cautious but compelling conclusion. Seasonal variation in cardiovascular events and physiology is well documented, yet direct evidence for seasonal variation in imaging-defined cardiac structure and function remains sparse, especially outside athlete cohorts. This makes cardiac MRI an important test case.

If a seasonal signal is detected in MRI-derived phenotypes, it should be interpreted as potentially biologically meaningful but likely modest in magnitude. Any observed effect will need careful discrimination from acquisition-year trends, scanner effects, body-size differences, and operational scheduling bias. Conversely, a null result would also be informative, because it would suggest that well-known cardiovascular seasonality in outcomes and hemodynamics does not necessarily translate into detectable variation in MRI-derived cardiac phenotype at the population level.

The current repository is intentionally transparent about access limitations. Several open-access full texts have been downloaded and summarized, while a number of important papers remain represented by verified bibliographic metadata and abstract-level extraction only. The review can therefore already guide analysis design, but it would benefit from a later pass with institutional full-text access before submission as a formal standalone review manuscript.
